% Section 2: Related Work (strengthened literature review)

\section{Related Work}
\label{sec:related}

We organize prior art along four streams that frame GBSS planning for BVLOS: (i)~multi-sensor C-UAS / UTM sensing, (ii)~surveillance network \emph{placement} for UAM/AAM, (iii)~multi-objective / metaheuristic placement methods, and (iv)~urban LoS geometry. Table~\ref{tab:related_comparison} summarizes closest systems work against this paper's dual-layer SCOPAS framing.

\begin{table*}[t]
\centering
\caption{Systems comparison of closely related GBSS / C-UAS / AAM surveillance placement work and method analogs. Tokens summarize scope; ``---'' means not the paper's focus.}
\label{tab:related_comparison}
\setlength{\tabcolsep}{3pt}
\begin{tabular}{lccccccc}
\toprule
\textbf{Work} & \textbf{Sensors} & \textbf{Urban geometry / LoS} & \textbf{Objectives} & \textbf{Optimizer} & \textbf{Coop$\neq$Noncoop} & \textbf{City-scale case} \\
\midrule
Guvenc et al.\ \cite{guvenc2018detection} & Multi (survey) & --- & Detection/tracking & --- & Partial & --- \\
Besada et al.\ \cite{besada2022cuas} & Radar/RF/EO/Ac & Sim.\ models & C-UAS for UTM & --- & Yes (noncollab.) & --- \\
Yan et al.\ \cite{yan2023passive} & RF/Ac/vision & Urban survey & Detection/tracking & --- & Implicit & --- \\
MUSCAT \cite{muscat2023} & Heterogeneous & DTED / elev. & $M_c$, $M_g$, cost & Exhaustive & No (fused) & Small / synthetic \\
Aievola et al.\ \cite{aievola2022ground} & Radar & Urban UAM & Coverage / radar net & MIP / analysis & No & Corridor-scale \\
Dulia \& Shihab (SAND) \cite{dulia2024sand} & Multi-type AAM & Terrain classes & Cost + $P_D$ floors & ILP / design & Discussed & Multi-city OH \\
Molina et al.\ \cite{molina2008layout} & Homog.\ WSN & Flat field & Nodes + energy & NSGA-II / MOEA & --- & Synthetic field \\
Iqbal et al.\ \cite{iqbal2015wsn} & WSN (survey) & --- & Coverage/cost/\ldots & MOEA survey & --- & --- \\
Charlish et al.\ \cite{charlish2016radar} & Radar & --- & Sensor management & GA & --- & --- \\
\textbf{This work} & Radar+RF+EO+Ac & OSM + ray-trace & $M_{wp,\mathrm{coop}}$, $M_{wp,\mathrm{noncoop}}$, CapEx & NSGA-II & \textbf{Yes} & City + SJC + Paulista lite \\
\bottomrule
\end{tabular}
\end{table*}


\subsection{Multi-Sensor C-UAS Sensing for UTM}

Amateur and small-UAV detection surveys emphasize complementary modalities---radar, RF/Remote~ID, EO/IR, and acoustic---and the value of fusion under urban clutter \cite{guvenc2018detection,yan2023passive}. Multisensor track fusion can maintain tracks under partial sensor outages \cite{jovanoska2018multisensor,dudczyk2022multisensor}, while small-RCS phenomenology remains challenging for radar alone \cite{semkin2020analyzing}. For UTM specifically, Besada et al.\ review Counter-UAS sensors and provide simulation models for noncollaborative surveillance feeding tactical UTM services \cite{besada2022cuas}. Blasch et al.\ argue that information fusion is an autonomy enabler for UTM \cite{blasch2021information}, consistent with classical JDL-style fusion handbooks \cite{liggins2009handbook}.

These works establish \emph{why} heterogeneous GBSS matter and that cooperative (e.g., RID/RF) and non-cooperative (radar/EO/acoustic) sensing address different threat classes. They typically stop at detection algorithms, sensor models, or system architecture---not CapEx-constrained placement under city-scale LoS occlusion with separate coop/noncoop coverage objectives.

\subsection{Surveillance Network Placement for UAM / AAM}

Closest placement literature includes MUSCAT, which formulates multi-sensor C-UAS placement with coverage/gap/cost metrics and exhaustive search over candidate sites, adopting elevation-oriented terrain models \cite{muscat2023}. Exhaustive enumeration guarantees small-$n$ optima but becomes intractable beyond roughly 10--15 sites. Aievola et al.\ study ground-based radar networks for UAM with design and performance analysis under urban constraints, leaning on mathematical programming / structured radar-network design rather than multi-modal CapEx Pareto fronts \cite{aievola2022ground}. Dulia and Shihab's SAND model designs AAM surveillance networks that minimize sensor cost subject to coverage and detection-probability floors over terrain classes, and discusses cooperative vs.\ non-cooperative sensing options in a cost--benefit setting \cite{dulia2024sand}. Resource-allocation MIP surveys underline scalability limits once formulations become non-convex \cite{katoh1998resource}.

Relative to these, our work keeps MUSCAT's metric lineage and heterogeneous mix, but (i)~optimizes threat-weighted \emph{dual-layer} $M_{wp,\mathrm{coop}}$ and $M_{wp,\mathrm{noncoop}}$ jointly with CapEx via NSGA-II, (ii)~uses OSM building geometry with deterministic ray-traced LoS, and (iii)~reports city/airport/Paulista fronts plus matched CapEx ablations (Controls~B--C) rather than a single fused score or exhaustive optimum claim.

\subsection{Multi-Objective and Metaheuristic Placement}

Holland's GA foundations \cite{holland1992adaptation} and Deb's NSGA-II \cite{deb2002fast} underpin Pareto search for discrete placement. WSN layout has long used MOEAs to trade coverage, node count, and energy \cite{molina2008layout,iqbal2015wsn,zou2019WSN}; radar sensor management has likewise used GAs \cite{charlish2016radar}. We treat these as \emph{method transfer}: coverage--cost Pareto search is mature in WSN/radar management, but transferring it to urban C-UAS GBSS requires occlusion-aware $P_D$ models and an explicit coop/noncoop objective split that WSN energy/connectivity formulations do not provide.

\subsection{Urban Geometry and Line-of-Sight Modeling}

Statistical urban channel models (e.g., 3GPP TR~38.901 LoS probability) and ITU short-range outdoor recommendations support radio planning without full building footprints \cite{3gpp38901,itur_p1411}. OpenStreetMap provides globally available building footprints used across geospatial and telecom studies \cite{osm_applications,elias2023exploring}. Deterministic ray-tracing on voxelized OSM buildings is more conservative than statistical LoS and yields altitude-stratified shadow zones---the geometry regime we adopt for GBSS planning.

\subsection{Gap Analysis}

Table~\ref{tab:related_comparison} highlights three persistent gaps: (1)~placement papers rarely expose \emph{separate} cooperative and non-cooperative coverage objectives under CapEx; (2)~exhaustive / MIP approaches struggle at city-scale candidate sets with heterogeneous types and occlusion physics; (3)~urban C-UAS placement validations seldom combine OSM geometry, ray-traced LoS, and multi-city (or multi-scene) Pareto evidence. This paper addresses (1)--(3) within a $g\cdot p$ NSGA-II budget, with Controls~A--C bounding small-$n$ quality and CapEx-matched modality claims---without asserting city-scale optimality gaps.
